% =====================================================================
\section*{How this work was done}
\addcontentsline{toc}{section}{How this work was done}
% =====================================================================

The authors chose the question, directed the work and judged every result. AI models from two developers derived, checked and wrote the mathematics. The theorem and its proof were derived by OpenAI's \OPENAIMODEL{}. Three agents of the same developer each checked the complete written proof, in separate runs that did not write it. Anthropic's Claude Opus~5.5 then re-derived the proof refutation-first, in two rounds: an independent derivation by hand, followed by an adversarial attack on its verdicts by a separate run that did not write them. This review found no false step and recommended three changes of presentation, which this manuscript makes: the passage from the quartered encoding to finite rigidity (\cref{prop:quarter-ti}(b)), the lemma that justifies conditioning on the background (\cref{lem:freeze}), and the orientation of the seed of the row markers (rule (Q2) and \cref{bg:markers}(c)). The search of the prior literature was done by Claude Opus~5.5. The manuscript was written by Claude Opus~5.5 from the audited proofs; this includes the explicit rules of \cref{sec:interaction} and the program of \cref{app:machine}, which the proofs specified in words. A separate run of OpenAI's \OPENAIMODEL{}, which had not seen the work, then reviewed the finished manuscript refutation-first, including the explicit rules, the tables and the program, and compared it with the original derivation. It checked finite parts with pairs of independently written programs: the scaffold tables and their seams, the program of $M$ on bounded inputs, and the arithmetic of the windows. It found no false step in the main results, and five local errors and five points of presentation or credit, which this manuscript corrects. Models are not listed as authors; the authors take responsibility for every claim.

Version~1.1.0 adds a citation of the authors' preprint \cite{DM26thermal}, which version~1.0.0 did not cite: two sentences comparing it with Theorem~A in the prior work of \cref{sec:intro}, a row of \cref{tab:landscape} and one sentence in \cref{sec:open}. Claude Opus~5.5 read that preprint and wrote these additions. No statement or proof has changed.

Version~1.2.0 adds the Lean formalization of \cref{app:lean} and three corrections of wording. The Lean proofs, their port to the toolchain pinned for release and the checks described in \cref{app:lean} were done by agents running Claude Opus~5.5, directed by a lead agent running the same model, which ruled on every choice of the formalization. The formalization found no false step. It found three imprecise sentences, corrected in this version: the kernel of a specification by factors was not specified where its normalization is $0$ (\cref{sec:setting}); the list of rules that involve payload coordinates omitted rule (K2) and the matching of colours (\cref{int:sub-split} and the proof of \cref{pay:factor}); and the list of sensor coordinates in \cref{sec:payload} included two components of colours that are $\circ$. No other statement or proof has changed. Claude Opus~5.5 wrote this version.

% =====================================================================
\section*{Verification status}
\addcontentsline{toc}{section}{Verification status}
% =====================================================================

\begin{itemize}
\item \textbf{Proved in the text.} Theorem~A (\cref{thm:main} with \cref{prf:windows}), Corollary~B (\cref{cor:undecidable}) and every lemma and proposition of \cref{sec:scaffold,sec:background,sec:payload,sec:proof} and of \cref{app:scaffold,app:machine}, given the inputs below.
\item \textbf{Imported with stated scope.} Existence and translation invariance of the plus and minus states of the Ising model on $\Z^2$, through the FKG inequality \cite{FriedliVelenik17}; Dobrushin's uniqueness theorem \cite{Dobrushin68,Georgii11}; and standard measure theory (regular conditional probabilities, compactness of the configuration space).
\item \textbf{Conditional.} \Cref{prop:exact} assumes Input~(I) of \cref{sec:exact}, the classical exact facts about the two-dimensional Ising model. Theorem~A and Corollary~B do not use it.
\item \textbf{The finite tables.} The interaction is given by explicit rules on a finite product of finite sets (\cref{sec:interaction}, \cref{app:tables}); the machine $M$ is given by its program (\cref{app:machine}), and its transition table contains the transition table of $U$. The tables were not listed in full. The rules were written for this manuscript from the specification of the audited proofs. Their agreement with that specification was checked by the writer; then by hand and refutation-first in a separate two-round review of this manuscript by Claude Opus~5.5, which derived every new rule and the program again and then attacked its own verdicts; and then in the final review by \OPENAIMODEL{}. No program listed or checked the whole allowed set, which depends on $U$.
\item \textbf{Checked by machine.} No proof in the text uses a computation. In the final review, independently written programs checked finite parts of the construction: the scaffold tables and their seams, the program of $M$ on bounded inputs, and the arithmetic of the windows. These checks support the text; they are not part of any proof.
\item \textbf{Formalized.} Theorem~A, Corollary~B and the existence of a machine $U$ as in \cref{set:sub-tm} are proved in Lean~4 with Mathlib, with no unproved step and only the three standard axioms (\cref{app:lean}). Theorem~A and Corollary~B are proved from two published results that Mathlib lacks, taken as hypotheses: Dobrushin's uniqueness theorem and the existence of the plus state of the two-dimensional Ising model, in the form of \cite{FriedliVelenik17}. \Cref{sec:exact}, \cref{prf:rem-midpoint,prf:rem-flip} and the $\Sigma^0_1$-completeness of $\mathcal H_U$ are not formalized. A comparator check confirms that the proved statements are those of a challenge file that imports only Mathlib, and two kernel checkers re-check the proofs. The Lean sources are in \texttt{lean/} of the repository.
\item \textbf{Checked less thoroughly.} The formal statements of \cref{app:lean} were written, and compared with the printed ones, by models of one developer only.
\item \textbf{Internal only.} Every review is internal; none of it is journal peer review.
\end{itemize}
